\documentclass[10pt,a4paper]{amsart}
\usepackage{amsmath,amssymb,amsthm,graphicx}
\usepackage{verbatim}
\newenvironment{reference}{\comment}{\endcomment}
%\newenvironment{reference}{}{}
\newenvironment{slogan}{\comment}{\endcomment}
\newenvironment{history}{\comment}{\endcomment}

% For commutative diagrams we use Xy-pic
\usepackage[all]{xy}

% We use 2cell for 2-commutative diagrams.
\xyoption{2cell}
% Theorem environments.
%
\theoremstyle{plain}
\newtheorem{theorem}[subsection]{Theorem}
\newtheorem{proposition}[subsection]{Proposition}
\newtheorem{lemma}[subsection]{Lemma}

\theoremstyle{definition}
\newtheorem{definition}[subsection]{Definition}
\newtheorem{example}[subsection]{Example}
\newtheorem{exercise}[subsection]{Exercise}
\newtheorem{situation}[subsection]{Situation}

\theoremstyle{remark}
\newtheorem{remark}[subsection]{Remark}
\newtheorem{remarks}[subsection]{Remarks}

\numberwithin{equation}{subsection}

% Macros
%
\def\lim{\mathop{\mathrm{lim}}\nolimits}
\def\colim{\mathop{\mathrm{colim}}\nolimits}
\def\Spec{\mathop{\mathrm{Spec}}}
\def\Hom{\mathop{\mathrm{Hom}}\nolimits}
\def\Ext{\mathop{\mathrm{Ext}}\nolimits}
\def\SheafHom{\mathop{\mathcal{H}\!\mathit{om}}\nolimits}
\def\SheafExt{\mathop{\mathcal{E}\!\mathit{xt}}\nolimits}
\def\Sch{\mathit{Sch}}
\def\Mor{\mathop{\mathrm{Mor}}\nolimits}
\def\Ob{\mathop{\mathrm{Ob}}\nolimits}
\def\Sh{\mathop{\mathit{Sh}}\nolimits}
\def\NL{\mathop{N\!L}\nolimits}
\def\CH{\mathop{\mathrm{CH}}\nolimits}
\def\proetale{{pro\text{-}\acute{e}tale}}
\def\etale{{\acute{e}tale}}
\def\QCoh{\mathit{QCoh}}
\def\Ker{\mathop{\mathrm{Ker}}}
\def\Im{\mathop{\mathrm{Im}}}
\def\Coker{\mathop{\mathrm{Coker}}}
\def\Coim{\mathop{\mathrm{Coim}}}

% Boxtimes
%
\DeclareMathSymbol{\boxtimes}{\mathbin}{AMSa}{"02}

%
% Macros for moduli stacks/spaces
%
\def\QCohstack{\mathcal{QC}\!\mathit{oh}}
\def\Cohstack{\mathcal{C}\!\mathit{oh}}
\def\Spacesstack{\mathcal{S}\!\mathit{paces}}
\def\Quotfunctor{\mathrm{Quot}}
\def\Hilbfunctor{\mathrm{Hilb}}
\def\Curvesstack{\mathcal{C}\!\mathit{urves}}
\def\Polarizedstack{\mathcal{P}\!\mathit{olarized}}
\def\Complexesstack{\mathcal{C}\!\mathit{omplexes}}
% \Pic is the operator that assigns to X its picard group, usage \Pic(X)
% \Picardstack_{X/B} denotes the Picard stack of X over B
% \Picardfunctor_{X/B} denotes the Picard functor of X over B
\def\Pic{\mathop{\mathrm{Pic}}\nolimits}
\def\Picardstack{\mathcal{P}\!\mathit{ic}}
\def\Picardfunctor{\mathrm{Pic}}
\def\Deformationcategory{\mathcal{D}\!\mathit{ef}}
\usepackage{cleveref}
\crefname{originalreference}{original reference}{original references}
\Crefname{originalreference}{Original reference}{Original references}
\makeatletter
\newlabel{remark-extension-by-zero}{{0G4N}{source}{}{sourcepointer-1}{}}
\newlabel{remark-extension-by-zero@cref}{{[originalreference][0][]0G4N}{[1][0][]source}{}{}{}}
\makeatother
\crefname{originalsection}{original section}{original sections}
\Crefname{originalsection}{Original section}{Original sections}
\makeatletter
\newlabel{flat-section-compactify}{{0ATT}{source}{}{sourcepointer-2}{}}
\newlabel{flat-section-compactify@cref}{{[originalsection][0][]0ATT}{[1][0][]source}{}{}{}}
\makeatother
\crefname{originallemma}{original lemma}{original lemmas}
\Crefname{originallemma}{Original lemma}{Original lemmas}
\makeatletter
\newlabel{lemma-well-defined-pre}{{0G4W}{source}{}{sourcepointer-3}{}}
\newlabel{lemma-well-defined-pre@cref}{{[originallemma][0][]0G4W}{[1][0][]source}{}{}{}}
\makeatother
\crefname{originallemma}{original lemma}{original lemmas}
\Crefname{originallemma}{Original lemma}{Original lemmas}
\makeatletter
\newlabel{flat-lemma-compactifications-cofiltered}{{0ATU}{source}{}{sourcepointer-4}{}}
\newlabel{flat-lemma-compactifications-cofiltered@cref}{{[originallemma][0][]0ATU}{[1][0][]source}{}{}{}}
\makeatother
\title{Extension by zero is compatible with pushing forward along a compactification}
\author{The Stacks Project Authors}
\date{stacks-project:duality.tex:local-snapshot-20261012; GNU FDL 1.2 or later}
\begin{document}
\maketitle
\noindent\textit{Editorial note.} All delivered mathematical text is the author's own, copied from the pinned Stacks Project chapter file duality.tex; no statement or proof step was rewritten, reordered or removed.\par\medskip
\noindent\textit{Editorial note.} Preamble provenance: the pinned chapter file opens with an input of the shared preamble and carries no preamble text of its own, so the project's single Stacks-wide preamble staged at sources/stacks-duality/preamble.tex is excerpted verbatim under the recipe fixed by LANE\_BRIEF 22.1, as three spans, lines 8-13 (verbatim package and the reference/slogan/history environments), lines 14-18 (Xy-pic setup) and lines 154-217 (the native theorem declarations and mathematical macros). Nothing was added, removed or redefined; the wrapper class and the amsmath/amssymb/amsthm/graphicx base packages are supplied by the assembler.\par\medskip
\noindent\textit{Editorial note.} The counted statement and proof contain no cross reference, no citation and no included file, so the delivered body is self contained.\par\medskip
\noindent\textit{Editorial note.} Canonical location of the counted proposition: Stacks tag 0G4X (sources/stacks-duality/pinned-original-source.tex line 8577).\par\medskip
\noindent\textit{Allowed standard prerequisites (editorial).}\begin{itemize}
\item The notion of a compactification over the base. More on Flatness, Section labelled flat-section-compactify, cited at proof lines 8612-8616. Fixes the meaning of the hypothesis.
\item The saturated case of the statement. Same chapter, Lemma labelled lemma-well-defined-pre, cited at proof lines 8618-8620. Settles the case j'(U) = f\textasciicircum{}\{-1\}(j(U)).
\item Cofilteredness of compactifications and its part (c). More on Flatness, Lemma labelled flat-lemma-compactifications-cofiltered part (c), cited at proof lines 8626-8630. Gives the two isomorphisms used in the general case.
\end{itemize}\medskip
\ifdefined\hypertarget\hypertarget{sourcepointer-1}{}\fi
\noindent\textit{Source reference.} Original reference 0G4N: sources/stacks-duality/pinned-original-source.tex, source lines 8203--8203. Same-chapter pointer into the author environment remark at source line 8203; invoked at proof line None of the counted proof. It is retained only as a navigation pointer; every load-bearing step is inside the delivered span.\par\medskip
\ifdefined\hypertarget\hypertarget{sourcepointer-2}{}\fi
\noindent\textit{Source reference.} Original section 0ATT: sources/stacks-chow/tags.txt, source lines 13626--13626. Cross-chapter standard prerequisite located in the project tag list staged in the local snapshot; invoked at proof line 8591 of the counted proof. It is retained only as a navigation pointer; every load-bearing step is inside the delivered span.\par\medskip
\ifdefined\hypertarget\hypertarget{sourcepointer-3}{}\fi
\noindent\textit{Source reference.} Original lemma 0G4W: sources/stacks-duality/pinned-original-source.tex, source lines 8472--8472. Same-chapter pointer into the author environment lemma at source line 8472; invoked at proof line 8593 of the counted proof. It is retained only as a navigation pointer; every load-bearing step is inside the delivered span.\par\medskip
\ifdefined\hypertarget\hypertarget{sourcepointer-4}{}\fi
\noindent\textit{Source reference.} Original lemma 0ATU: sources/stacks-chow/tags.txt, source lines 13627--13627. Cross-chapter standard prerequisite located in the project tag list staged in the local snapshot; invoked at proof line 8606 of the counted proof. It is retained only as a navigation pointer; every load-bearing step is inside the delivered span.\par\medskip
\begin{lemma}
\label{lemma-well-defined}
Let $j : U \to X$ be an open immersion of Noetherian schemes.
Let $j' : U \to X'$ be a compactification of $U$ over $X$ (see proof)
and denote $f : X' \to X$ the structure morphism.
Then we have a canonical isomorphism $Rj_! \to Rf_* \circ R(j')_!$
of functors $D^b_{\textit{Coh}}(\mathcal{O}_U) \to
\text{Pro-}D^b_{\textit{Coh}}(\mathcal{O}_X)$ where
$Rj_!$ and $Rj'_!$ are as in Remark \ref{remark-extension-by-zero}.
\end{lemma}

\begin{proof}
The fact that $X'$ is a compactification of $U$ over $X$ means
precisely that $f : X' \to X$ is proper, that $j'$ is an open immersion,
and $j = f \circ j'$. See
More on Flatness, Section \ref{flat-section-compactify}.
If $j'(U) = f^{-1}(j(U))$, then the lemma follows immediately from
Lemma \ref{lemma-well-defined-pre}.
If $j'(U) \not = f^{-1}(j(U))$, then denote $X'' \subset X'$ the
scheme theoretic closure of $j' : U \to X'$ and denote
$j'' : U \to X''$ the corresponding open immersion.
Picture
$$
\xymatrix{
& & X'' \ar[d]^{f'} \\
& & X' \ar[d]^f \\
U \ar[rr]^j \ar[rru]^{j'} \ar[rruu]^{j''} & & X
}
$$
By
More on Flatness, Lemma \ref{flat-lemma-compactifications-cofiltered} part (c)
and the discussion above we have isomorphisms
$Rf'_* \circ Rj''_! = Rj'_!$ and $R(f \circ f')_* \circ Rj''_! = Rj_!$.
Since $R(f \circ f')_* = Rf_* \circ Rf'_*$ we conclude.
\end{proof}

\end{document}
